\section{Related Work}
\noindent{\textbf{Multi-Modal Representation Learning.}} In zero-shot 3D classification, recent research has focused on multi-modal learning with limited 3D data by distilling CLIP knowledge~\cite{qi2023contrast, hegde2023clip, huang2023clip2point, xue2024ulip, liu2024openshape} because CLIP possesses extensive knowledge in vision-language tasks. For example, MixCon3D~\cite{gao2024sculpting} improved the performance of zero-shot 3D classification by contrastive learning across the image-text-point cloud using CLIP knowledge. 
In addition, ULIP-2~\cite{xue2024ulip} is a multimodal pre-training that automatically generates comprehensive language descriptions for 3D shapes without manual annotation.
These conventional methods achieve high zero-shot performance on ModelNet40~\cite{wu20153d}, ScanObjectNN~\cite{uy2019revisiting}, and Objaverse-LVIS~\cite{deitke2023objaverse}. However, these methods basically focus only on improving training methods, and there are limits to how much performance can be improved with this approach alone. We believe it is essential for the model and data sides to evolve together. This paper thus explores the potential of synthetic dataset expansion to achieve scalable geometric learning for zero-shot 3D classification.

\noindent{\textbf{Text-to-3D Models.}} Recent text-to-image models have experienced rapid growth. Inspired by this, text-to-3D models have also become a prominent area of research. 
% Recent text-to-image models allow us to generate realistic images, and these achievement techniques extended to 3D vision. 
Text-to-3D models also face the challenge of limited 3D data, making it difficult to generate open-vocabulary 3D data directly from text. To address this, recent text-to-3D models employ a strategy that leverages text-to-2D models or CLIP's extensive language and image knowledge. They first generate images or image features and then lift these images to 3D data~\cite{poole2022dreamfusion, lin2023magic3d, chen2023fantasia3d, wang2024prolificdreamer, tsalicoglou2024textmesh}.
% In text-to-3D models, the mainstream trend is to combine CLIP with NeRF~\cite{mildenhall2021nerf} and Gaussian Splatting~\cite{kerbl20233d}, as there is still the constraint of using a small amount of 3D data~\cite{poole2022dreamfusion, lin2023magic3d, chen2023fantasia3d, wang2024prolificdreamer, tsalicoglou2024textmesh}.
For example, DreamFusion~\cite{poole2022dreamfusion} efficiently generates 3D data using CLIP's knowledge and NeRF's gradient updates.
In addition, Zero123-XL~\cite{liu2023zero} trains on the relatively large 3D dataset called Objaverse-XL~\cite{deitke2024objaverse} and achieves highly accurate 3D generation. Recently, it has also been applied to explicit 3D representations such as point clouds and meshes~\cite{nichol2022point, jun2023shap}. Using a text-to-image and image-to-point cloud pipeline, Point-E~\cite{nichol2022point} quickly generates point clouds that correspond to text prompts.
In this paper, we utilize text-to-3D models for dataset expansion.
% When using text-to-3D models for dataset expansion, the generation speed is a critical factor. Therefore, we use Point-E as the text-to-3D model, as it is both fast and capable of producing consistent quality.

\noindent{\textbf{Generative Models as Data++.}} 
Recent generative models have evolved to generate realistic synthetic data that is visually almost indistinguishable~\cite{rombach2022high, saharia2022photorealistic, nichol2021glide}. And then the data generated from generative models can be thought of as data for recognition models with controllability and rich representation as data++~\cite{isola2022when, sariyildiz2023fake}. It reduces the collecting cost with real-world data while enabling the efficient generation of meaningful synthetic datasets. For example, StableRep~\cite{tian2024stablerep} learns visual representations by incorporating synthetic images generated by Stable Diffusion~\cite{rombach2022high} into multi-positive contrastive learning.  
StableRep is designed to learn images comprehensively across different views by treating multiple synthetic images of the same input text prompt as positive samples. 
In addition, SynCLR~\cite{tian2024learning} performs contrastive learning using only synthetic images and captions.
Despite training without real images, it performs equally well and better than traditional self-supervised learning such as DINOv2~\cite{oquab2023dinov2} and OpenCLIP~\cite{cherti2023reproducible} in image classification and semantic segmentation. These studies have shown that it is possible to build high performance recognition models while training only synthetic data and keeping data collection costs low by learning visual representations.
Inspired by these studies, we utilize data generated by generative models for contrastive learning.
